%% ma: L10-B02-0024
%% lop: 10
%% chuong: 1
%% bai: 2
%% dang: 10.2.5
%% loai: TLN
%% muc: TH
%% dap_an: "5"
%% nguon: "Ảnh giáo viên gửi 10/10/2026 (ThS. Nguyễn Hoàng Việt, Chương 2 Tập hợp), A-Đề 1, Phần 3 Câu 20"
%% kien_thuc: "Bài 2 (tập hợp và các phép toán trên tập hợp)"
%% trang_thai: chinh-thuc
%% ghi_chu: "Đề gốc yêu cầu hai số (1 và 4); theo yêu cầu giáo viên đã sửa thành tổng số phần tử của A giao B và A hiệu B"
\begin{ex}
Cho hai tập hợp $A=\{n\in\mathbb N\mid-2<n\le4\}$ và $B=\{x\in\mathbb Z\mid2x^3+x^2-x=0\}$. Tính tổng số phần tử của hai tập hợp $A\cap B$ và $A\setminus B$.
\shortans{5}
\loigiai{$A=\{0;1;2;3;4\}$. $2x^3+x^2-x=x(2x-1)(x+1)=0\Leftrightarrow x\in\left\{0;\dfrac12;-1\right\}$, nên $B=\{-1;0\}$. $A\cap B=\{0\}$ có $1$ phần tử; $A\setminus B=\{1;2;3;4\}$ có $4$ phần tử. Tổng bằng $1+4=5$.}
\end{ex}
